\section{Backend Architecture and Implementation}

The backend of the DBaaS platform serves as the centralized orchestration engine and traffic gateway, fundamentally architected to support highly available, multi-tenant database provisioning at scale. It is engineered entirely in Go (Golang)—a strategic technological selection necessitated by Go's highly efficient concurrency primitives (goroutines and channels) and its position as the lingua franca of the cloud-native ecosystem. This implementation achieves deterministic, low-latency HTTP request processing and highly performant TCP stream multiplexing. Furthermore, it facilitates seamless, native interoperability with the underlying container orchestration layer via the \texttt{client-go} SDK.

To enforce rigorous fault-domain isolation and permit independent horizontal scaling, the system architecture implements a strict decoupling of the Control Plane from the Data Plane, physically separating operations into two distinct microservices:

\begin{enumerate}
    \item \textbf{Control Plane (REST API Server):} A stateless, API-first HTTP gateway responsible for declarative infrastructure orchestration, distributed session management, and automated API generation. It securely manages the lifecycle of Kubernetes Custom Resource Definitions (CRDs) and maintains the global platform state without intercepting persistent data traffic.
    \item \textbf{Data Plane (PGProxy SNI Router):} A highly specialized Layer 4 TCP proxy engineered to dynamically route raw PostgreSQL wire-protocol traffic to isolated Kubernetes pods via Transport Layer Security Server Name Indication (TLS SNI). By operating strictly below the Application Layer (Layer 7), it circumvents the HTTP API infrastructure entirely.
\end{enumerate}

This deliberate architectural bifurcation ensures that computationally intensive, long-running analytical queries executed within the data plane cannot induce resource starvation within the control plane, thereby yielding an exceptionally resilient and fault-tolerant distributed system.

\subsection{Control Plane: REST API Architecture}

The REST API Server functions as the orchestration engine and control plane for the DBaaS platform. To ensure high maintainability, testability, and strict separation of concerns, the codebase strictly adheres to the principles of Clean Architecture and Domain-Driven Design (DDD). The system is highly modularized via interface-driven development and constructor-based Dependency Injection, organized into three distinct tiers:

\begin{itemize}
    \item \textbf{Handler Layer (Transport \& Routing):} Implemented utilizing the highly performant \texttt{Gin} HTTP framework, this layer acts as the system's edge. It intercepts incoming RESTful requests, enforces rate-limiting middleware, performs strict schema validation on Data Transfer Objects (DTOs), and returns standardized, deterministic JSON API responses. Crucially, handlers contain no business logic; they securely delegate validated payloads to the underlying services.
    \item \textbf{Service Layer (Domain Logic):} Acting as the core of the application, the service layer encapsulates the platform's proprietary business rules. It orchestrates complex, multi-step distributed operations such as OAuth2 authentication state machines, unified database backups, and dynamic asynchronous infrastructure provisioning. Specialized domain abstractions are implemented for distinct paradigms, such as relational PostgreSQL schema evolution and unstructured MongoDB document manipulation.
    \item \textbf{Repository Layer (Persistence Abstraction):} This layer abstracts all I/O operations and state persistence, completely decoupling the domain logic from the underlying storage engines. It implements the Repository Pattern to manage interactions with the centralized Platform Meta-Database via \texttt{pgxpool} (for ACID-compliant metadata storage) and Redis (for ephemeral session states), effectively isolating raw SQL logic from the rest of the application.
\end{itemize}

\subsection{Multi-Tenant Kubernetes Provisioning}

Rather than interacting with low-level container runtimes directly, the backend is architected as an intelligent Kubernetes client. The internal \textbf{Operator Provisioner} module interfaces natively with the Kubernetes (K3s) API server to execute declarative, desired-state infrastructure management:

\begin{itemize}
    \item \textbf{Operator Pattern Orchestration:} Upon receiving a project creation request, the provisioner dynamically templates and applies Custom Resource Definitions (CRDs) for \texttt{CloudNativePG} (for PostgreSQL) or \texttt{MongoDBCommunity} (for MongoDB). By delegating cluster lifecycle management to specialized K8s Operators, the system natively inherits advanced capabilities such as automated pod scheduling, self-healing, streaming replication, and automated failover without requiring custom backend intervention.
    \item \textbf{Namespace Isolation and RBAC:} To enforce strict multi-tenant security boundaries, every database project is provisioned into a dynamically generated, cryptographically unique Kubernetes Namespace (e.g., \texttt{pg-<uuid>}). The control plane interacts with these namespaces using a least-privilege \texttt{ServiceAccount}, bound by rigorous Role-Based Access Control (RBAC) policies. These policies restrict the backend's permissions strictly to the localized manipulation of CRDs and Secrets, mitigating the blast radius of any potential compromise.
    \item \textbf{Cloud-Agnostic Portability:} By abstracting the database lifecycle entirely through standard Kubernetes operators rather than relying on proprietary managed services (e.g., AWS RDS or Google Cloud SQL), the DBaaS platform achieves true cloud-agnosticism. The entire control plane and data plane can be seamlessly deployed across any public cloud provider, private on-premise infrastructure, or edge environment running Kubernetes, fundamentally eliminating vendor lock-in.
\end{itemize}

\subsection{Dynamic Database Connection Management}

A fundamental challenge in multi-tenant database architectures is efficiently routing and executing user queries against thousands of isolated database instances without exhausting file descriptors or memory. The platform resolves this via an intelligent, thread-safe \textbf{Connection Manager}:

\begin{itemize}
    \item \textbf{Lazy Initialization and JIT Resolution:} Rather than maintaining an exhaustively large, static connection pool to all active databases, the Connection Manager instantiates tenant-specific multiplexed connection pools (\texttt{pgxpool} for PostgreSQL, native \texttt{mongo-driver} for MongoDB) entirely just-in-time (JIT). When an API request targets a specific project, the manager dynamically resolves the corresponding internal Kubernetes service endpoints and injects the required credentials on the fly.
    \item \textbf{Pool Caching and Eviction Lifecycle:} Once initialized, the connection pool is securely cached in memory to eliminate TLS handshake and authentication overhead for subsequent requests. To prevent memory leaks and resource starvation, the system implements an aggressive eviction policy: idle pools are dynamically pruned, and active pools are forcefully closed and evicted from memory when their corresponding database instance is paused, scaled to zero, or deleted.
\end{itemize}

\subsection{Data Plane: PGProxy SNI Router}

Facilitating direct, secure SQL access for external GUI clients (e.g., DBeaver, DataGrip) and CLI tools (\texttt{psql}) presents a complex routing challenge in a dynamically orchestrated Kubernetes environment. Exposing individual NodePorts or LoadBalancers per tenant is architecturally prohibitive. The platform solves this via a proprietary \textbf{PGProxy SNI Router}, a custom Layer 4 TCP proxy:

\begin{itemize}
    \item \textbf{TLS Server Name Indication (SNI) Routing:} External clients initiate connections to the proxy via deterministically generated subdomains (e.g., \texttt{db-<uuid>.postgres.domain.com}). The proxy implements zero-copy TLS peeking to parse the \texttt{ClientHello} packet, extracting the requested SNI extension to dynamically resolve the target tenant's internal cluster IP.
    \item \textbf{Zero-Termination End-to-End Encryption:} Unlike traditional reverse proxies, PGProxy strictly operates at the transport layer and does \textit{not} terminate TLS. By establishing a transparent, bidirectional TCP tunnel directly to the internal CloudNativePG pod, the system guarantees cryptographically secure end-to-end encryption. This architectural decision crucially preserves native PostgreSQL authentication mechanisms, such as SCRAM-SHA-256 channel binding, which are otherwise broken by standard Man-in-the-Middle (MITM) termination.
\end{itemize}

\subsection{Dynamic Auto-API Generation (PostgREST)}

To eliminate the need for users to write boilerplate backend middleware, the DBaaS platform natively provisions a fully featured, schema-driven REST API for every PostgreSQL tenant. This is accomplished via the automated orchestration of \textbf{PostgREST} sidecar deployments:

\begin{itemize}
    \item \textbf{Automated Role-Based Authorization:} Upon database initialization, the operator provisioner natively connects to the cluster as a superuser to establish a strict, principle-of-least-privilege RBAC topology. It configures an \texttt{authenticator} role for connection multiplexing, an \texttt{anon} role for unauthenticated access, and an \texttt{app\_user} role granted scoped CRUD privileges over the public schema.
    \item \textbf{Dynamic Ingress and Cryptographic Security:} The backend cryptographically generates a 64-byte JWT secret and pre-signs a stateless service-role API key during the provisioning phase. It subsequently deploys a Kubernetes \texttt{Service} and binds it to a dynamically templated Traefik HTTP \texttt{IngressRoute} (e.g., \texttt{rest-<uuid>.domain.com}). This enables users to perform RESTful interactions over an automatically generated, highly secure gateway instantaneously.
\end{itemize}

\subsection{Multi-Model Database Integration}

To support diverse workloads ranging from strict ACID-compliant transactions to flexible document storage, the platform abstracts the complexities of multi-paradigm database management behind a unified, polymorphic REST API. Requests are dynamically routed to specialized service implementations based on the underlying persistence engine:

\begin{itemize}
    \item \textbf{Relational Persistence (PostgreSQL):} Exposes a robust API for deterministic schema evolution, dynamic DDL execution, ACID-compliant row manipulations, and complex relational aggregates.
    \item \textbf{NoSQL Document Storage (MongoDB):} Facilitates the ingestion and querying of unstructured BSON documents, dynamic collection management, and the execution of highly parallelized aggregation pipelines.
    \item \textbf{Unified Disaster Recovery Pipeline:} A highly generalized backup abstraction layer dispatches format-specific logical export tasks. It initiates asynchronous streams of \texttt{pg\_dump} archives or \texttt{mongodump} BSON payloads, piping the binary streams directly over HTTP to the client. This achieves a cohesive, zero-disk-footprint disaster recovery interface irrespective of the underlying storage model.
\end{itemize}

\subsection{Platform Meta-Database Architecture}

While tenant data is rigidly isolated across independent Kubernetes clusters, the global state and configuration of the DBaaS platform are centrally managed within an ACID-compliant \textbf{Meta-Database} (a dedicated PostgreSQL instance). The schema design strictly prioritizes cryptographic security and high-throughput metadata retrieval:

\begin{itemize}
    \item \textbf{Identity and Access Management (IAM):} The \texttt{users} schema manages global identity, utilizing highly secure hashing algorithms for credential storage, while tracking temporal session metadata and role-based privilege assignments.
    \item \textbf{Project State Registry:} The \texttt{projects} schema serves as the canonical source of truth for all provisioned infrastructure. It maintains referential integrity to the IAM schema while tracking polymorphic database types, hierarchical resource tiers, and state-machine transitions (e.g., \texttt{creating}, \texttt{running}, \texttt{failed}).
    \item \textbf{Zero-Trust Credential Security:} In adherence to zero-trust security paradigms, tenant database credentials and Data Source Names (DSNs) are \textit{never} persisted within the Meta-Database. To neutralize the threat of lateral movement during a data breach, the Go backend directly encapsulates the tenant passwords into Base64-encoded Kubernetes \texttt{Secret} resources prior to cluster creation. The backend securely retrieves these secrets on-the-fly via the \texttt{InstanceDSNService} to route queries, ensuring passwords exist only in ephemeral memory.
\end{itemize}

\subsection{Authentication \& Security Architecture}

Securing both the control plane API and the tenant instances is a primary architectural focus. The backend implements a robust, stateless authentication system:

\begin{itemize}
    \item \textbf{OAuth and Identity Federation:} The platform supports seamless Google OAuth integration alongside standard email/password registration. User passwords are cryptographically hashed using industry-standard algorithms before persistence.
    \item \textbf{JWT Token State Machine:} Upon successful authentication, the API issues short-lived Access Tokens (e.g., 15 minutes) and long-lived Refresh Tokens (e.g., 30 days). Access Tokens are entirely stateless, enabling extremely fast, decentralized authorization across the microservices without requiring synchronous database lookups.
    \item \textbf{Distributed Session Management:} To maintain strict control over active sessions and enable instant revocation, Refresh Tokens are centrally persisted in an external \textbf{Redis} cluster. During the token refresh cycle, the backend rigorously validates the token against Redis and completely rotates the token pair, effectively neutralizing token hijacking attacks.
\end{itemize}



\subsection{System Workflows: Lifecycle of a DBaaS Project}

The architectural sophistication of the platform is best illustrated by tracing the declarative, event-driven lifecycle of a tenant database provisioning request:

\begin{enumerate}
    \item \textbf{Request Interception and Semantic Validation:} The API Server intercepts the incoming payload via the \texttt{ProjectService}, executing strict semantic validation against the requested persistence engine, cryptographic password complexity, and resource-tier quotas.
    \item \textbf{Optimistic State Registration:} Upon validation, a canonical record is synchronously committed to the Meta-Database with an initial \texttt{status} of \texttt{creating}. This optimistic state update ensures a non-blocking HTTP response, enabling the frontend to immediately transition to a reactive loading state.
    \item \textbf{Asynchronous Declarative Orchestration:} The computationally intensive provisioning phase is decoupled into a background goroutine. The \texttt{OperatorProvisioner} securely injects the tenant password into a Kubernetes \texttt{Secret}, dynamically templates the required Custom Resource definitions (e.g., a CloudNativePG \texttt{Cluster}), and asynchronously applies them against the K3s API server.
    \item \textbf{Operator-Driven Reconciliation:} The specialized Kubernetes controllers natively detect the CRD submission and begin their reconciliation loops. They autonomously orchestrate the underlying pod scheduling, mount highly available Persistent Volume Claims (PVCs), configure internal ClusterIP networking, and securely mount the injected credentials to bootstrap the daemon.
    \item \textbf{State Fulfillment and JIT Connection:} The backend continuously polls the cluster for readiness. Upon successful reconciliation, the Meta-Database state is transitioned to \texttt{running}. Subsequent query executions trigger the \texttt{ConnectionManager} to retrieve the injected K8s \texttt{Secret}, dynamically construct the DSN, multiplex the connection pool, and execute the payload against the isolated tenant.
    \item \textbf{Infrastructure Teardown and Eviction:} During project deletion, the backend initiates a cascading teardown sequence. It synchronously removes the Meta-Database registry, evicts any active connection pools from RAM to prevent dangling references, and deletes the associated CRD, commanding the Kubernetes operator to systematically terminate all underlying tenant infrastructure.
\end{enumerate}
